\chapter{Time Stepping: Accuracy and Stability}
\label{chpt:time}

% ---------------------------------------------------------------------------
% Placeholder for figures that have not been made yet. Replace each
% \TSfigplaceholder{...} with \includegraphics once the figure exists.
% ---------------------------------------------------------------------------
\providecommand{\TSfigplaceholder}[1]{%
  \fbox{\parbox[c][#1][c]{0.85\textwidth}{\centering\sffamily\small
  \textbf{[Figure to be added]}}}}

Every problem we have seen so far has been \emph{steady}, which means the unknown $\phi$ does not change with time. In this chapter, we will see how to handle problems where $\phi$ changes in time using
\emph{time stepping} or \emph{time integration}. There are many time stepping methods, and each comes with its own costs. Two questions guide our comparison of them:
\begin{enumerate}
  \item \textbf{Accuracy.} How fast does the error shrink as the time step $\Delta t$ is made smaller?
  \item \textbf{Stability.} Does the numerical solution stay bounded at all, or can it blow up even when the true solution decays?
\end{enumerate}
Answering both lets us weigh the cost of each method against what it gives us.
Our model problem will be the heat equation, the parabolic member of the three
classic PDEs introduced in Chapter~\ref{chpt:intro2Classs}.
\index{time stepping}

\section{The Heat Equation}
\label{sec:ts-heat}

Let $\Omega = [a,b]\times[c,d]\subset\mathbb{R}^2$ and
let $\phi(x,y,t)$ be temperature. We solve for $\phi$ satisfying
\begin{equation}
  \frac{\partial \phi}{\partial t} = k\,\nabla^2\phi + f(x,y,t),
  \qquad (x,y)\in\Omega,\quad t > 0,
  \label{eq:ts-heat}
\end{equation}
where $k>0$ is the diffusivity (i.e., how quickly the heat spreads) and $f$ is a known forcing function (i.e., a heat source/sink). Unlike the
steady problems of earlier chapters, the solution requires
\begin{itemize}
  \item an \textbf{initial condition}, $\phi(x,y,0) = \phi_0(x,y)$ for all $(x,y)\in\Omega$, and
  \item \textbf{boundary conditions} on $\partial\Omega$ for every $t>0$, exactly as in
        Section~\ref{sec:2Dbcs}, except the boundary data may now depend on time.
\end{itemize}
\index{heat equation}\index{diffusion equation}\index{initial condition}

Equation~\eqref{eq:ts-heat} is also called the \emph{diffusion equation}, because
temperature is only one example. The quantity $\phi$ can be many things: temperature (heat), the concentration of a
dissolved chemical, the density of a pollutant in the air, or even the density
of people wandering randomly around a room. In every case, the Laplacian moves
$\phi$ from where there is more of it to where there is less, so the heat
equation \emph{smooths things out}. Sharp features are rounded off, peaks
flatten, and in some cases the solution relaxes toward a steady state.

\subsection{A Random Walk}
\label{sec:ts-walk}
\index{random walk}

Why does diffusion look like a Laplacian? Consider a person who has had a bit
too much to drink and is walking along a line, starting from the point $x_0$.
Every $\Delta t$ seconds they stumble a distance $\Delta x$, either to the right
with probability $p = 0.5$, or to the left with probability $1-p = 0.5$. A single walker's
path is random and jagged, as seen in the left panel of Figure~\ref{fig:ts-walk}. However, if we start $N_w$ walkers at
the same point $x_0$ and let them all stumble independently, we can describe the
\emph{distribution} of their positions at a given time. Let $P(t,x)$ be the
probability of finding a walker at position $x$ at time $t$. To be at $x$ at
time $t + \Delta t$, a walker must have been at $x - \Delta x$ and stumbled
right, or at $x + \Delta x$ and stumbled left, so
\begin{equation}
  P(t + \Delta t, x) = \tfrac12\,P(t, x - \Delta x) + \tfrac12\,P(t, x + \Delta x).
  \label{eq:ts-master}
\end{equation}
Taylor expanding both sides of \eqref{eq:ts-master}, the way we did for finite
differences in Section~\ref{sec:secondderiv}, the left-hand side expands in time,
\begin{equation*}
  P(t + \Delta t, x) = P + \Delta t\,\frac{\partial P}{\partial t}
  + \frac{\Delta t^2}{2}\,\frac{\partial^2 P}{\partial t^2}
  + \mathcal{O}(\Delta t^3),
\end{equation*}
and the two terms on the right-hand side expand in space,
\begin{align*}
  P(t, x + \Delta x) &= P + \Delta x\,\frac{\partial P}{\partial x}
  + \frac{\Delta x^2}{2}\,\frac{\partial^2 P}{\partial x^2}
  + \frac{\Delta x^3}{6}\,\frac{\partial^3 P}{\partial x^3}
  + \frac{\Delta x^4}{24}\,\frac{\partial^4 P}{\partial x^4}
  + \mathcal{O}(\Delta x^5), \\
  P(t, x - \Delta x) &= P - \Delta x\,\frac{\partial P}{\partial x}
  + \frac{\Delta x^2}{2}\,\frac{\partial^2 P}{\partial x^2}
  - \frac{\Delta x^3}{6}\,\frac{\partial^3 P}{\partial x^3}
  + \frac{\Delta x^4}{24}\,\frac{\partial^4 P}{\partial x^4}
  + \mathcal{O}(\Delta x^5),
\end{align*}
where $P$ and its derivatives on the right are all evaluated at $(t,x)$. Averaging
the two spatial expansions, the odd derivatives cancel, just as for the central
difference \eqref{eq:fdCentral}:
\begin{equation*}
  \tfrac12\,P(t, x - \Delta x) + \tfrac12\,P(t, x + \Delta x)
  = P + \frac{\Delta x^2}{2}\,\frac{\partial^2 P}{\partial x^2}
  + \frac{\Delta x^4}{24}\,\frac{\partial^4 P}{\partial x^4}
  + \mathcal{O}(\Delta x^6).
\end{equation*}
Setting the two sides equal, cancelling $P$, and dividing by $\Delta t$ gives
\begin{equation}
  \frac{\partial P}{\partial t}
  = \frac{\Delta x^2}{2\,\Delta t}\,\frac{\partial^2 P}{\partial x^2}
  \underbrace{- \frac{\Delta t}{2}\,\frac{\partial^2 P}{\partial t^2}
  + \frac{\Delta x^4}{24\,\Delta t}\,\frac{\partial^4 P}{\partial x^4}
  + \cdots}_{\text{higher-order terms}} = \frac{\Delta x^2}{2\,\Delta t}\,\frac{\partial^2 P}{\partial x^2}
  + \text{(higher-order terms)}.
  \label{eq:ts-walkheat}
\end{equation}
Apart from the higher-order terms, this is the heat equation in one dimension, with diffusivity
$k = \Delta x^2 / (2\Delta t)$.

\begin{figure}[htbp]
  \centering
  \TSfigplaceholder{5cm}
  \caption{Left: the position of a few random walkers plotted against time, each taking steps of $\pm\Delta x$ every $\Delta t$, with vertical lines marking $t = 1$ and $t = 10$. Right: the distribution $P(t,x)$ of many walkers' positions (as histograms, i.e.\ data) at $t=1$ and $t=10$, with the corresponding solutions of the heat equation \eqref{eq:ts-walkheat} drawn as curves. The later distribution should be visibly wider and shorter.}
  \label{fig:ts-walk}
\end{figure}

The higher-order terms of \eqref{eq:ts-walkheat} are negligible when $P$ barely changes over a single step. When does that happen? We can think of our drunk walkers are molecules moving randomly. At the macroscopic level, the amount the individual molecules change with space ($\Delta x$) with respect to time ($\Delta t$) are very small. So
let's consider \eqref{eq:ts-walkheat} as $\Delta x \to 0$ and $\Delta t \to 0$.
\begin{itemize}
  \item If $\Delta t \propto \Delta x$, the coefficient behaves like $\Delta x \to 0$
        and the limit is $\partial P/\partial t = 0$: the walkers never spread.
  \item If $\Delta t \propto \Delta x^2$, the coefficient stays fixed, and choosing
        $\Delta t = \Delta x^2/(2k)$ gives the heat equation with diffusivity $k$.
  \item If $\Delta t \propto \Delta x^3$, the coefficient behaves like
        $1/\Delta x \to \infty$: the walkers spread infinitely far in any
        positive time.
\end{itemize}
Only $\Delta t \propto \Delta x^2$ gives a finite, nonzero limit. At this scaling, the higher-order terms of \eqref{eq:ts-walkheat} go to zero. 

\begin{remark}
Keep the scaling $\Delta t \sim \Delta x^2$ in mind; it will reappear in
Section~\ref{sec:ts-stability}.
\end{remark}
\index{diffusive scaling}


\begin{figure}[htbp]
  \centering
  \TSfigplaceholder{4.5cm}
  \caption{A probability distribution $P(t,x)$ obtained from many random walkers, plotted at three times (e.g.\ $t=0$, $t\approx 1.3$, $t\approx 11$). The earliest distribution is narrow and noisy; later ones are progressively wider, flatter, and smoother. Each time should use a distinct color \emph{and} line style.}
  \label{fig:ts-smoothing}
\end{figure}

\section{Exact Solution}
\label{sec:ts-mol}

With the random walk as motivation, we return to the two-dimensional heat
equation \eqref{eq:ts-heat}. We already know how to spatially discretize the Laplacian $\nabla^2$. Let's consider this problem on a Type~I grid with Dirichlet boundary conditions on the domain $\Omega = [0,\ell]^2$ such that for $N$ points,
\begin{equation}
  x_i = i\,\Delta x, \quad y_j = j\,\Delta x, \qquad i,j = 1,\dots,N,
  \qquad \Delta x = \frac{\ell}{N+1}.
  \label{eq:ts-grid}
\end{equation}
We can discretize \eqref{eq:ts-heat} in terms of the Laplacian $L$ and flatten the unknown field $\boldsymbol{\phi}$ to obtain a large system of ordinary differential equations (ODEs),
\begin{equation}
  \frac{d\boldsymbol{\phi}}{dt} = k\,L\,\boldsymbol{\phi} + \mathbf{g}(t),
  \qquad \boldsymbol{\phi}(0) = \boldsymbol{\phi}_0,
  \label{eq:ts-ode}
\end{equation}
where $\mathbf{g}(t)$ collects the forcing $f$ evaluated at the grid points
and the boundary contributions (which now may change in time). 

\begin{remark}
Note we have separated the problem into two independent choices: how
to discretize space and how to step in time. The spatial error and the
temporal error then add together, just as the discretization and solver errors did in
Section~\ref{sec:gmres-tol}.
\end{remark}

Before discretizing time, it helps to look at the exact solution of
\eqref{eq:ts-ode}. Let $k=1$ and $\mathbf{g}=\mathbf{0}$. For a scalar
ODE $y' = \lambda y$, we would write $y(t) = e^{\lambda t}y(0)$. The exact solution to \eqref{eq:ts-ode} is then
\begin{equation}
  \boldsymbol{\phi}(t) = e^{tL}\,\boldsymbol{\phi}_0,
  \qquad
  e^{tL} = I + tL + \frac{t^2}{2}L^2 + \frac{t^3}{3!}L^3 + \cdots,
  \label{eq:ts-matexp}
\end{equation}
where the matrix exponential is defined by its Taylor series, which converges far too slowly for this to be a practical solution. Since $L$ is symmetric, it can be decomposed as
\begin{equation}
  L = P\,D\,P^{\mathsf T},
  \qquad D = \operatorname{diag}(\lambda_1,\lambda_2,\dots,\lambda_n),
  \label{eq:ts-diag}
\end{equation}
where the columns of $P$ are the eigenvectors and $\lambda_m$ the eigenvalues of $L$. Let $\mathbf{y} = P^{\mathsf T}\boldsymbol{\phi}$. Multiplying
$\boldsymbol{\phi}' = kL\boldsymbol{\phi}$ by $P^{\mathsf T}$ gives
$\mathbf{y}' = kD\mathbf{y}$. Each row of this system represents the scalar equation
$$y_m' = k\lambda_m\,y_m$$. 
Defining $\hat{\lambda}_m = k \lambda_m$, we can consider one of these scalar equations as
\begin{equation}
  \frac{dy_m}{dt} = \hat{\lambda}_m\,y_m, \qquad y_m(0) = y_m^0,
  \qquad\text{with exact solution}\qquad y_m(t) = e^{\hat{\lambda}_m t}\,y_m^0.
  \label{eq:ts-test}
\end{equation}

\begin{remark}
If we understand what a time stepper does to \eqref{eq:ts-test} for one
eigenvalue $\lambda_m$, we understand what it does to the whole system, one
eigenvalue at a time. Each of the following time steppers, applied to
$\boldsymbol{\phi}' = kL\boldsymbol{\phi}$, is the same as applying it to every $y_m$
separately.
\end{remark}

This reduces the question to what the biggest $\lambda$ is. For the system we have defined thus far, the eigenvalues are
\begin{equation}
  \lambda_{p,q} = -\frac{4}{\Delta x^2}\left[
  \sin^2\!\left(\frac{p\pi\Delta x}{2\ell}\right) +
  \sin^2\!\left(\frac{q\pi\Delta x}{2\ell}\right)\right],
  \qquad p,q = 1,\dots,N.
  \label{eq:ts-eigs}
\end{equation}
\index{Laplacian!eigenvalues}
Every eigenvalue is real and negative, so every mode of the true solution decays.
The eigenvalue of largest magnitude belongs to the most oscillatory mode
$p = q = N$ (a checkerboard), and since the $\sin^2$ terms are then nearly one,
\begin{equation}
  \lambda_{\max} \approx -\frac{8}{\Delta x^2}.
  \label{eq:ts-lammax}
\end{equation}
We have thus defined the largest possible eigenvalue, which we help us with our stability analysis for each temporal scheme.


\section{Forward and Backward Euler}
\label{sec:ts-euler}
\index{forward Euler}\index{backward Euler}

Discretization of time is very similar to that for space. Consider the time interval $[0,T]$ evaluated over $N_t$ equally spaced points, such that
$\Delta t = T/N_t$. Then, we can define time levels, or evaluations of the system at specific times, as
\begin{equation}
  t^n = n\,\Delta t, \qquad n = 0,1,\dots,N_t,
  \qquad \boldsymbol{\phi}^n \approx \boldsymbol{\phi}(t^n).
  \label{eq:ts-levels}
\end{equation}
Superscripts denote the time level, while subscripts still denote the
grid point, so $\phi^n_{i,j} \approx \phi(x_i, y_j, t^n)$. The simplest approximation of the
time derivative is the first-order difference
\[
  \frac{d\boldsymbol{\phi}}{dt} \approx \frac{\boldsymbol{\phi}^{n+1} - \boldsymbol{\phi}^n}{\Delta t},
\]
and identical to Section~\ref{sec:firstderiv}), we will define forward and backwards approximations of this derivative.

\paragraph{Forward Euler.} Evaluate $\mathbf{F}$ at the \emph{known} level $t^n$:
\begin{equation}
  \frac{\boldsymbol{\phi}^{n+1} - \boldsymbol{\phi}^n}{\Delta t} = \mathbf{F}(t^n,\boldsymbol{\phi}^n)
  \quad\Longrightarrow\quad
  \boxed{\boldsymbol{\phi}^{n+1} = \left(I + \Delta t\,k L\right)\boldsymbol{\phi}^n + \Delta t\,\mathbf{g}(t^n).}
  \label{eq:ts-FE}
\end{equation}
Each step is a sparse matrix-vector multiplication. 
\index{explicit method}

\paragraph{Backward Euler.} Evaluate $\mathbf{F}$ at the \emph{unknown} level $t^{n+1}$:
\begin{equation}
  \frac{\boldsymbol{\phi}^{n+1} - \boldsymbol{\phi}^n}{\Delta t} = \mathbf{F}(t^{n+1},\boldsymbol{\phi}^{n+1})
  \quad\Longrightarrow\quad
  \boxed{\left(I - \Delta t\,k L\right)\boldsymbol{\phi}^{n+1} = \boldsymbol{\phi}^n + \Delta t\,\mathbf{g}(t^{n+1}).}
  \label{eq:ts-BE}
\end{equation}
Now each step requires a linear \emph{solve}. A solve is going to be more expensive than a multiplication, so there must be a tradeoff that makes backward Euler useful. Finding that tradeoff is the goal of the next two sections.
\index{implicit method}

\begin{remark}
A method in which the new value is given
directly by known values is called \emph{explicit}. A method whose new value is
defined only by unknown values is called \emph{implicit}.
\end{remark}


\begin{remark}
Starting from $\boldsymbol{\phi}^0 = \boldsymbol{\phi}_0$, it takes exactly
$N_t$ steps to reach $\boldsymbol{\phi}^{N_t} \approx \boldsymbol{\phi}(T)$. A loop that runs one step too many or too
few compares your solution at $T \pm \Delta t$ with the exact solution at $T$, an error that looks like a first-order
method no matter what you implemented. Likewise, if $\Delta t$ does not divide $T$
evenly, you can never land on $T$. Compute the time of each level as $t^n = n\Delta t$
rather than repeatedly adding $\Delta t$, which accumulates round-off.
\end{remark}


\section{Accuracy}
\label{sec:ts-accuracy}
\index{order of accuracy}

To compare the two methods, apply both to the test equation \eqref{eq:ts-test}. Forward Euler gives
$y^{n+1} = (1 + \Delta t\lambda)\,y^n$ and backward Euler gives
$(1-\Delta t\lambda)\,y^{n+1} = y^n$. Repeating either $n$ times, we can
write the \emph{exact} solutions of the two approximations,
\begin{align}
  \text{exact:} \qquad & y(t^n) = e^{\lambda t^n}\,y^0, \label{eq:ts-exactscalar}\\
  \text{FE:}    \qquad & y^n = \left(1 + \Delta t\,\lambda\right)^n y^0, \label{eq:ts-FEscalar}\\
  \text{BE:}    \qquad & y^n = \left(\frac{1}{1 - \Delta t\,\lambda}\right)^n y^0. \label{eq:ts-BEscalar}
\end{align}
The error of forward Euler at $t^n$, measured relative to the size of the initial data, is
\begin{equation}
  \frac{|y^n - y(t^n)|}{|y^0|} = \left|\left(1 + \Delta t\lambda\right)^n - e^{\lambda t^n}\right|
  = \frac12\,\lambda^2\,t^n\,e^{\lambda t^n}\,\Delta t + \mathcal{O}(\Delta t^2).
  \label{eq:ts-FEerror}
\end{equation}
We will see in the next section what can go wrong.
The error is proportional to $\Delta t$, so forward Euler is first-order
accurate.


\begin{remark}
The test equation gives us exact formulas, but for a general method we use the
method of manufactured solutions.
\end{remark}

\begin{definition}[Local truncation error]
The \emph{local truncation error} $\tau^n$ of a time stepper is the error accumulated from the iterative scheme starting from an exact solution per unit of time. \end{definition}
\index{local truncation error}\index{global error}

For forward Euler, $\tau^n = \frac{\boldsymbol{\phi}(t^{n+1})-\boldsymbol{\phi}(t^n)}{\Delta t} -
\mathbf{F}(t^n,\boldsymbol{\phi}(t^n)) = \frac{\Delta t}{2}\boldsymbol{\phi}''(t^n) + \cdots$,
which is exactly the forward difference error \eqref{eq:fdFO1}. 



The
local errors made at each of the $N_t = T/\Delta t$ steps accumulate into \textit{global error}. If $\tau^n = \mathcal{O}(\Delta t^p)$ and
the method is \textbf{stable}, then the \textit{global error} at a fixed time $T$ is also
$\mathcal{O}(\Delta t^p)$ and the method is said to be \emph{$p$-th order accurate}. If the method is unstable, global error becomes unbounded and your system will not converge.

\subsection{Measuring the Order of Accuracy}
\label{sec:ts-measure}
\index{convergence study}

In practice, we check the order of a time stepper numerically. We use a manufactured solution that now also depends on time, compute the solution at a fixed
final time $T$ for a sequence of time steps, and measure the relative error
\begin{equation}
  \text{Err}(\Delta t) =
  \frac{\norm{\boldsymbol{\phi}^{N_t} - \boldsymbol{\phi}_{\text{exact}}(T)}}
       {\norm{\boldsymbol{\phi}_{\text{exact}}(T)}}.
  \label{eq:ts-relerr}
\end{equation}


\begin{figure}[htbp]
  \centering
  \TSfigplaceholder{6cm}
  \caption{A convergence study for backward Euler on the unforced test problem $\phi = e^{-2kt}\sin(x)\sin(y)$ on $[0,2\pi]^2$. Relative error \eqref{eq:ts-relerr} at a fixed $T$ versus $\Delta t$ on log-log axes for two spatial resolutions $N$, drawn with distinct markers and colors, with a dashed $\mathcal{O}(\Delta t)$ reference line. For small $\Delta t$ each curve should flatten at a level set by the spatial error, with the finer grid flattening lower.}
  \label{fig:ts-convergence}
\end{figure}



\section{Stability}
\label{sec:ts-stability}
\index{stability}

Look again at the error \eqref{eq:ts-FEerror} and the test equation with
$\lambda < 0$, which we know is the case for every eigenvalue of the
Laplacian. The exact solution $e^{\lambda t}y^0$ decays to zero as $t\to\infty$, so we expect the numerical approximation to as well.

Consider the forward Euler case, we only see a decay in the solution when 
\begin{equation}
  \left|1 + \Delta t\,\lambda\right| < 1.
  \label{eq:ts-FEstab}
\end{equation}
This is called a \emph{stability restriction} and it must hold for \emph{every}
eigenvalue. Recalling that we determined the maximum eigenvalue as $\lambda = k\lambda_{\max} \approx -8k/\Delta x^2$ in \eqref{eq:ts-lammax}, 

\begin{align*}
  -1 &< 1 + \Delta t\left(\frac{-8k}{\Delta x^2}\right) < 1 \\
  -2 &< \Delta t\left(\frac{-8k}{\Delta x^2}\right) < 0 \\
  \frac14 &> \frac{k\,\Delta t}{\Delta x^2} > 0,
\end{align*}
so forward Euler for the two-dimensional heat equation requires
\begin{equation}
  \boxed{\frac{4\,k\,\Delta t}{\Delta x^2} < 1.}
  \label{eq:ts-FElimit}
\end{equation}
If this relation holds, the error \eqref{eq:ts-FEerror} decays along with the
solution for every mode. If it does not, the most oscillatory mode is multiplied
by a number less than $-1$ every step, flips sign, grows, and soon swamps
the solution with a growing checkerboard pattern. Notice the scaling: $\Delta t$ must shrink like $\Delta x^2$, the
diffusive scaling of Section~\ref{sec:ts-walk}.
\index{stability restriction}

\begin{figure}[htbp]
  \centering
  \TSfigplaceholder{5.5cm}
  \caption{Forward Euler applied to the test equation $y' = \lambda y$ with $\lambda<0$. The exact solution (curve) is compared with forward Euler (symbols) for three time steps: $|1+\Delta t\lambda|$ well below one (smooth decay), $-1 < 1+\Delta t\lambda < 0$ (decay with alternating sign), and $1+\Delta t\lambda < -1$ (oscillating growth). Backward Euler at the largest time step may be added for comparison.}
  \label{fig:ts-decay}
\end{figure}

\paragraph{Backward Euler.} Repeating the argument with \eqref{eq:ts-BEscalar}, backward Euler decays only if
\begin{equation}
  \left|\frac{1}{1 - \Delta t\,\lambda}\right| < 1.
  \label{eq:ts-BEstab}
\end{equation}
For $\lambda<0$ and any $\Delta t > 0$, the denominator $1 - \Delta t\lambda$ is
larger than one, so \eqref{eq:ts-BEstab} always holds. We end up with
\emph{no} restriction on $\Delta t$ at all, which we refer to as
\emph{unconditionally stable}.
\index{unconditional stability}


\subsection{The Tradeoff Between Forward and Backward Euler}
\label{sec:ts-tradeoff}

Recall that backward Euler requires a linear solve, making it more expensive than forward. There is a tradeoff between the two, however, and it is related to their stabilities.
\begin{itemize}
  \item Backward Euler is an \textbf{implicit} method. Each step needs a linear solve, but there are no (or few) restrictions on $\Delta t$ for stability.
  \item Forward Euler is an \textbf{explicit} method. There is no system to solve, so each step is very fast, but $\Delta t$ is restricted by \eqref{eq:ts-FElimit}.
\end{itemize}

\textit{Stability is not accuracy.} Backward Euler is stable for an enormous $\Delta t$, but
its answer will be inaccurate. The point is that with an implicit
method, $\Delta t$ can be chosen by the accuracy you need, while with an explicit
method $\Delta t$ may be forced far smaller than accuracy requires.




\section{Higher-Order Time Steppers}
\label{sec:ts-higher}

Stability tells us which time steps we may take, and accuracy tells us how good
the answer is. Both Euler methods are only first-order accurate. Just as the central difference of
Section~\ref{sec:central} improved on the one-sided differences, we can do better
in time. We describe four methods for the general system
$\boldsymbol{\phi}' = \mathbf{F}(t,\boldsymbol{\phi})$, where we have defined $\mathbf{F}(t,\boldsymbol{\phi}) = kL\boldsymbol{\phi} + \mathbf{g}(t)$ for the heat equation.

\subsection{Crank--Nicolson}
\label{sec:ts-CN}
\index{Crank--Nicolson}

Instead of evaluating $\mathbf{F}$ at one end of the step, average the two ends:
\begin{equation}
  \frac{\boldsymbol{\phi}^{n+1} - \boldsymbol{\phi}^n}{\Delta t}
  = \frac12\left[\mathbf{F}(t^n,\boldsymbol{\phi}^n) + \mathbf{F}(t^{n+1},\boldsymbol{\phi}^{n+1})\right].
  \label{eq:ts-CN}
\end{equation}
This is the trapezoidal rule, known as the \emph{Crank--Nicolson} method when applied to PDEs. The difference quotient is now
\emph{centered} at $t^{n+1/2}$, which is the source of its extra order of accuracy (recall \eqref{eq:fdCentral}!). Since $\boldsymbol{\phi}^{n+1}$ appears on the
right-hand side, it is implicit.


\subsection{Explicit Midpoint (RK2)}
\label{sec:ts-RK2}
\index{Runge--Kutta methods}\index{explicit midpoint method}

Another way to center the step at $t^{n+1/2}$ is to first take half a forward Euler step
to estimate the solution at the midpoint, and then use the slope there for the full step:
\begin{align}
  \boldsymbol{\phi}^{*} &= \boldsymbol{\phi}^n + \frac{\Delta t}{2}\,\mathbf{F}(t^n, \boldsymbol{\phi}^n),
  \label{eq:ts-RK2a}\\
  \boldsymbol{\phi}^{n+1} &= \boldsymbol{\phi}^n + \Delta t\,\mathbf{F}\!\left(t^n + \tfrac{\Delta t}{2},\, \boldsymbol{\phi}^{*}\right).
  \label{eq:ts-RK2b}
\end{align}
The intermediate value $\boldsymbol{\phi}^*$ is called a \emph{stage}. Methods built from several
stages are \emph{Runge--Kutta} (RK) methods. This one, the explicit midpoint method, is a
second-order RK method (RK2). It requires no solve, so it is explicit.


\subsection{BDF2}
\label{sec:ts-BDF2}
\index{BDF methods}\index{multistep methods}

Runge--Kutta methods use extra stages within one step. \emph{Multistep} methods instead
reuse the solution from earlier steps. The \emph{backward differentiation formulas} (BDF) approximate
$d\boldsymbol{\phi}/dt$ at the new level $t^{n+1}$ by a one-sided difference through
several past levels, and evaluate $\mathbf{F}$ at $t^{n+1}$, as backward Euler does
(backward Euler is BDF1). The second-order member is
\begin{equation}
  \frac{3\boldsymbol{\phi}^{n+1} - 4\boldsymbol{\phi}^{n} + \boldsymbol{\phi}^{n-1}}{2\,\Delta t}
  = \mathbf{F}(t^{n+1},\boldsymbol{\phi}^{n+1}).
  \label{eq:ts-BDF2}
\end{equation}

BDF2 cannot start on its own. To compute $\boldsymbol{\phi}^1$ it needs
$\boldsymbol{\phi}^{0}$ and $\boldsymbol{\phi}^{-1}$, and the latter does not exist. The first step
must be taken with a one-step method (or, when testing with a manufactured
solution, $\boldsymbol{\phi}^1$ can be taken from the exact solution).



\subsection{Adaptive Runge--Kutta (RK45)}
\label{sec:ts-RK45}
\index{adaptive time stepping}

Rather than fixing $\Delta t$, \emph{adaptive} methods choose it as they go. The popular
RK45 (Dormand--Prince) pair computes a fourth- and a fifth-order approximation from the
same stages. The difference of these approximations estimates the error of the step, and $\Delta t$ is grown or shrunk to keep that estimate
below a user-supplied tolerance. Writing one is tedious, so we use the
open implementations \texttt{ode45} in MATLAB and
\texttt{scipy.integrate.solve\_ivp(..., method='RK45')} in python. Both take a function
that evaluates $\mathbf{F}(t,\boldsymbol{\phi})$, the time interval, the initial condition, and
relative and absolute tolerances, and both can report how many times $\mathbf{F}$ was
evaluated.

\begin{table}[h!]
    \centering
    \begin{tabular}{|>{\raggedright\arraybackslash}p{3.6cm}|>{\raggedright\arraybackslash}p{2.2cm}|>{\raggedright\arraybackslash}p{2.2cm}||>{\raggedright\arraybackslash}p{4.2cm}|}
        \hline
        \textit{Method} & \textit{Type} & \textit{Order} & \textit{Stable for all $z<0$?} \\
        \hline
        Forward Euler \eqref{eq:ts-FE} & explicit & 1 & no, only $-2 \le z \le 0$ \\
        Backward Euler \eqref{eq:ts-BE} & implicit & 1 & yes \\
        Crank--Nicolson \eqref{eq:ts-CN} & implicit & 2 & yes \\
        Explicit midpoint \eqref{eq:ts-RK2b} & explicit & 2 & no, only $-2 \le z \le 0$ \\
        BDF2 \eqref{eq:ts-BDF2} & implicit & 2 & yes \\
        RK45 & explicit & 4/5 (adaptive) & no, only about $-3.3 \le z \le 0$ \\
        \hline
    \end{tabular}
    \caption{Summary of the time steppers in this chapter. Here $z = \Delta t\,\lambda$, where $\lambda$ is an eigenvalue of $kL$.}
    \label{tab:ts-summary}
\end{table}

\section{Solving the Implicit Systems Efficiently}
\label{sec:ts-solve}
\index{implicit method!linear solves}

Implicit methods buy their stability with a linear solve at every step, so the
cost of that solve deserves attention. Backward Euler, Crank--Nicolson, and BDF2 all reduce each time step to a linear
system of the form
\begin{equation}
  A\,\boldsymbol{\phi}^{n+1} = \mathbf{b}^n, \qquad n = 0,\dots,N_t-1,
  \label{eq:ts-Axb}
\end{equation}
where $\mathbf{b}^n$ changes every step but, as long as $\Delta t$ is fixed, the matrix
$A$ does \emph{not}. For backward Euler, $A = I - \Delta t\,kL$. We have to solve
\eqref{eq:ts-Axb} thousands of times, so it pays to think about how.

\paragraph{1) Use GMRES.} Solve each step with GMRES as in Chapter~\ref{chpt:GMRES},
using the previous time level as the initial guess (Section~\ref{sec:gmres-x0}),
\[
  \boldsymbol{\phi}^{n+1} = \texttt{gmres}(A,\ \mathbf{b}^n,\ \text{guess} = \boldsymbol{\phi}^{n}).
\]
If $\Delta t$ is small, $\boldsymbol{\phi}^n$ is already very close to $\boldsymbol{\phi}^{n+1}$ and only a few
iterations are needed. Recall from Section~\ref{sec:gmres-tol} that the tolerance should
sit just below the other errors in the problem.

\paragraph{2) Factor once, solve many times.} Since $A$ never changes, we can pay to
factor it once, before the time loop, and then reuse the factors at every step. If $A$
is \emph{symmetric positive definite} (SPD), it has a \emph{Cholesky factorization}
\begin{equation}
  A = R\,R^{\mathsf T},
  \label{eq:ts-chol}
\end{equation}
where $R$ is lower triangular and called the \emph{Cholesky factor}. Then each solve
$RR^{\mathsf T}\boldsymbol{\phi}^{n+1} = \mathbf{b}^n$ splits into two triangular solves,
\begin{equation}
  \left\{
  \begin{aligned}
    R\,\mathbf{z}^n &= \mathbf{b}^n && \text{($R$ lower triangular: forward substitution)},\\
    R^{\mathsf T}\boldsymbol{\phi}^{n+1} &= \mathbf{z}^n && \text{($R^{\mathsf T}$ upper triangular: back substitution)},
  \end{aligned}
  \right.
  \label{eq:ts-trisolve}
\end{equation}
each of which costs about as much as a multiplication by $R$.
\index{Cholesky factorization}

For the five-point Laplacian on an $N\times N$ grid with $n = N^2$ unknowns, a good factorization costs about $n^{3/2}$ operations and produces a factor with about $n\log n$ nonzeros,
so each pair of triangular solves costs about $n\log n$. The factorization is paid once; the solves
are paid every step.

\begin{remark}
The factor is only reusable while $A$ is unchanged. If $\Delta t$ changes, as in an
adaptive method, or $k$ depends on time, the factorization must be redone.
\end{remark}


\section{Numerical Implementation}
\label{sec:ts-code}

The code below sets up the two-dimensional heat equation on $[0,2\pi]^2$ with Dirichlet
conditions, using $N$ interior points per direction as in Chapter~\ref{chpt:GMRES}. It
builds the Laplacian, factors an implicit matrix once before the time loop, and lays out the
time loop with the bookkeeping of Section~\ref{sec:ts-euler}. The manufactured
solution, the forcing (with boundary contributions), and the update performed in each time step are left for you to write.
The backward Euler matrix is factored as an example; other implicit methods need their own matrix.

\subsection{MATLAB}
\lstinputlisting[language=matlab]{TimeStepping/code/heat_2d_timestep.m}

\subsection{python}
SciPy does not include a sparse Cholesky factorization, so instead we use
the sparse $LU$ factorization package \texttt{scipy.sparse.linalg.factorized} that returns a
function for repeated solves. Its cost behaves similarly for this problem.

\lstinputlisting[language=python]{TimeStepping/code/heat_2d_timestep.py}